\documentclass{sigchi}

% Use this command to override the default ACM copyright statement (e.g. for preprints). 
% Consult the conference website for the camera-ready copyright statement.


%% EXAMPLE BEGIN -- HOW TO OVERRIDE THE DEFAULT COPYRIGHT STRIP -- (July 22, 2013 - Paul Baumann)
% \toappear{Permission to make digital or hard copies of all or part of this work for personal or classroom use is 	granted without fee provided that copies are not made or distributed for profit or commercial advantage and that copies bear this notice and the full citation on the first page. Copyrights for components of this work owned by others than ACM must be honored. Abstracting with credit is permitted. To copy otherwise, or republish, to post on servers or to redistribute to lists, requires prior specific permission and/or a fee. Request permissions from permissions@acm.org. \\
% {\emph{CHI'14}}, April 26--May 1, 2014, Toronto, Canada. \\
% Copyright \copyright~2014 ACM ISBN/14/04...\$15.00. \\
% DOI string from ACM form confirmation}
%% EXAMPLE END -- HOW TO OVERRIDE THE DEFAULT COPYRIGHT STRIP -- (July 22, 2013 - Paul Baumann)


% Arabic page numbers for submission. 
% Remove this line to eliminate page numbers for the camera ready copy
\pagenumbering{arabic}


% Load basic packages
\usepackage{balance}  % to better equalize the last page
\usepackage{graphics} % for EPS, load graphicx instead
\usepackage{times}    % comment if you want LaTeX's default font
\usepackage{url}      % llt: nicely formatted URLs
\usepackage{color}

\usepackage{pdflscape}
\usepackage{longtable}
% \usepackage{geometry}
\usepackage[strict]{changepage}
\usepackage{german}


% llt: Define a global style for URLs, rather that the default one
\makeatletter
\def\url@leostyle{%
  \@ifundefined{selectfont}{\def\UrlFont{\sf}}{\def\UrlFont{\small\bf\ttfamily}}}
\makeatother
\urlstyle{leo}


% To make various LaTeX processors do the right thing with page size.
\def\pprw{210mm}
\def\pprh{297mm}
\special{papersize=\pprw,\pprh}
\setlength{\paperwidth}{\pprw}
\setlength{\paperheight}{\pprh}
\setlength{\pdfpagewidth}{\pprw}
\setlength{\pdfpageheight}{\pprh}

% Make sure hyperref comes last of your loaded packages, 
% to give it a fighting chance of not being over-written, 
% since its job is to redefine many LaTeX commands.
\usepackage[pdftex]{hyperref}
\hypersetup{
pdftitle={Supporting public deliberation through spatially enhanced dialogues},
pdfauthor={LaTeX},
bookmarksnumbered,
pdfstartview={FitH},
colorlinks,
citecolor=black,
filecolor=black,
linkcolor=black,
urlcolor=black,
breaklinks=true,
}

% create a shortcut to typeset table headings
\newcommand\tabhead[1]{\small\textbf{#1}}

% End of preamble. Here it comes the document.
\begin{document}

\title{Supporting public deliberation\\through spatially enhanced dialogues}
\subtitle{Master thesis}

\numberofauthors{1}
\author{
  \alignauthor Gerald Pape\\
    \affaddr{Institute for Geoinformatics}\\
    \email{g.pape@uni-muenster.de}\\
%   \alignauthor 2nd Author Name\\
%     \affaddr{Affiliation}\\
%     \affaddr{Address}\\
%     \email{e-mail address}\\
%     \affaddr{Optional phone number}    
%   \alignauthor 3rd Author Name\\
%     \affaddr{Affiliation}\\
%     \affaddr{Address}\\
%     \email{e-mail address}\\
%     \affaddr{Optional phone number}
}

\maketitle

\begin{abstract}
swaghetti yolonaise
\end{abstract}

% \keywords{
% 	Guides; instructions; author's kit; conference publications;
% 	keywords should be separated by a semi-colon.
% 	\textcolor{red}{Mandatory section to be included in your final version.}
% }

% \category{H.5.m.}{Information Interfaces and Presentation (e.g. HCI)}{Miscellaneous}

% See: \url{http://www.acm.org/about/class/1998/}
% for more information and the full list of ACM classifiers
% and descriptors. 
% \textcolor{red}{Mandatory section to be included in your
% final version. On the submission page only the classifiers'
% letter-number combination will need to be entered.}

\section{Introduction}
Since their first appearance, Web 2.0 applications utilized their collaborative character to gather information and opinions from their users. Today, modern information technologies are ubiquitous in many aspects of daily life. Involving citizens in decision processes around public matters through such applications has formed the field of ``eParticipation''. Its premise is to strengthen democratic processes between citizens and its governments through said modern information technologies \cite{Saebo_eParticipation, Medaglia2012_eParticipation}. One of many aspects is public deliberation which revolves around engaging citizen in dialogues about 


\section{Related Work}

\subsubsection{Argumentation mapping}
Rinner\cite{Rinner_ArgumentationMaps}\dots

Existing implementations\dots

Evaluation\dots
\subsubsection{Public deliberation and eParticipation}


\section{Approach}



\subsubsection{DialogMap}

In order to test the initial idea of supporting public deliberation through spatially enhanced dialogues, a working prototype had to be developed. Starting from an initial survey of existing research, a first prototypical application was developed. This prototype was then extended and refined with practical advice from members of a scientific citizens' initiative. Their input ranged from general suggestions to opinions of specific features. This chapter will give some details of design and implementation of the developed developed.


\subsubsection{Design decisions} % and features

As seen in X,Y and Z, important aspects of A are\dots

Internally, the prototype uses few data models. Contributions contain a title, description, two categories, a tags field, a favored counter, an optional time restriction field for start and ending times, an optional image, an optional reference to a parent contribution and optional references to child contributions. The parent and child contribution references create a simple parent-child connection between contributions, as children inherit the categories, tags, time restriction and title. A contribution serves both as a topic and as response to a topic. A contribution also contains references to features, references to features and references to URLs.\\
Features are geospatial entities with a spatial location, a reference to its contribution and properties for styling\footnote{\url{https://github.com/mapbox/simplestyle-spec}}.\\
Feature references contain a description of the featurea and the reference to a feature. URL references contain hyperlinks and a description of the hyperlink. The description of a contribution contains the text typed by a user with specially encoded references to features, URL references and feature references.\\
After signing in, users can create contributions in the manner of creating topics or writing responses to existing topics. Users have an e-mail address and a name.

The front page of the prototype puts a map side by side with a sidebar at right hand side containing from top to bottom the input form for new contributions, filter options, sorting order selector and a list of contributions. The input form consists of input fields for title, categories, time restriction, image and description. The description field allows the creation of spatial features and URL/feature references through connecting words with spatial representations or URLs.\\
A free text input field and multiple checkboxes allow to restrict the listed contribution as well as the geo-features displayed in the map. It is also possible to change the order of the list of contribution through a drop down field.\\
The list of contributions contains colored rectangles representing the different topics. Each box contains the title, time of writing, name of the author, categories, tags and the amount of times the contribution has been favored by users. It also contains a link which navigates to the replies written to the topic. A click on the contribution box expands it, revealing the description of the current topic.\\
After clicking the ``reply'' link, only the selected topic and replies are shown in the sidebar in a chronological order. In this view, each contribution shows the description by default as well as author and time and date of writing. The author of the contribution is able to edit the contribution. Other users are able to favor the contribution to show interest or agreement.\\
The map view contains a base map and several markers and polygons in different colors and different icons in case of markers. These relate to the contributions and are connected through the references in the description of the contributions. Which spatial features are displayed is determined through the state of the sidebar. In the topics overview, only the features created for the starting contribution are displayed in order to prevent cluttering of the view-port. When only the topic and its replies are displayed in the sidebar, all features related to the topic and its replies are shown on the map.\\
To emphasize the relationship between a contribution and its spatial features, a two way highlighting has been implemented. Hovering over either a contribution-box, marked word or spatial feature on the map triggers visual highlighting on all related contributions, marked words and spatial features. This allows to quickly grasp the relationship between features and contributions.\\
Users are able to use either traditional sing-up/sign-in methods or social sign-in through different providers to authenticate to the system.

% \begin{itemize}
% \item Map view with sidebar on right hand side
% \item Two way highlighting between contributions in sidebar and features on map
% \item Creation of Topic with \begin{itemize}
%     \item Title
%     \item Category/ies \begin{itemize}
%         \item in this specific case for two dimensions \begin{itemize}
%             \item Color
%             \item Icon
%             \end{itemize}
%         \end{itemize}
%     \item Tags
%     \item Time limit
%     \item Image
%     \item Special Description field which allows to create \begin{itemize}
%         \item Points and Polygons
%         \item References to existing Points and Polygons
%         \item Hyperlinks
%         \end{itemize}
%     \end{itemize}
% \item Sorting
% \item Filter \begin{itemize}
%     \item Fulltext
%     \item Categories
%     \item Tags
%     \item Time
% \end{itemize}
% \item Favorites
% \item Register/Sign in \begin{itemize}
%     \item with Google,Facebook, Twitter
%     \end{itemize}
% \end{itemize}


\subsubsection{Implementation}
\textit{DialogMap} has been implemented as a single-page web application using AngularJS\footnote{\url{http://angularjs.org/}} and Ruby on Rails\footnote{\url{http://rubyonrails.org/}}. The single-page structure was chosen in order to provide the user with a clear navigation between the overview and contribution answers. This also allows for a seamless browsing experience without full reloads of the page. AngularJS is a JavaScript framework with features like templating, two-way binding and DOM manipulation. It follows the model-view-controller pattern in order to bring server side paradigms to client-side development. AngularJS was chosen because of its popularity, extensibility and high number of available libraries. It also enables to wrap existing JavaScript libraries to be used in AngularJS context.\\
The mapping library Leaflet\footnote{\url{http://leafletjs.com/}} serves as base for displaying base maps and geospatial data. The user-facing web page was developed using tools like CoffeeScript\footnote{\url{http://coffeescript.org/}}, Haml\footnote{\url{http://haml.info/}} and Sass\footnote{\url{http://sass-lang.com/}} to speed up the development. The web page was developed with all major browsers in mind.\\
On the server side, components were developed using the Ruby on Rails framework with PostgreSQL\footnote{\url{http://www.postgresql.org/}}/PostGIS\footnote{\url{http://postgis.net/}} as data storage. Ruby on Rails, originally a full-stack model-view-controller web framework, is used as a JSON serving application logic. It was chosen because of its maturity and high number of available libraries. Front- and backend of the prototype communicate in REST\footnote{Representational State Transfer}-API\footnote{Application programming interface} like manner. This allows for easily replaceable front- and backend application stacks.\\
Figure \ref{fig:screenshot} shows the front page of the prototype with an active two way highlight.\\
Without the extensive use of open source software and code, development would have taken much longer. It is planned to release the source code through github\footnote{\url{https://github.com/ubergesundheit/dialogmap}}.

\begin{figure}[!h]
    \centering
    \includegraphics[width=0.9\columnwidth]{screenshot}
    \caption{Screenshot of the front page of \textit{DialogMap} with active highlight of a contribution and spatial feature.}
    \label{fig:screenshot}
\end{figure}


\section{Evaluation}



Interviews

Utility evaluation

Types of questions

Results

\section{Conclusion}

This work discusses the implementation and pre-evaluation of an prototype to support public deliberation through spatially enhanced dialogues.

\subsubsection{Future Work}
Pick up shortcomings emerged during evaluation. Point to solutions...

Legal implications of running such a website have to be explored.

% \section{Figure and Table Example}

% \begin{figure}[!h]
% \centering
% \includegraphics[width=0.9\columnwidth]{Figure1}
% \caption{With Caption Below, be sure to have a good resolution image
%   (see item D within the preparation instructions).}
% \label{fig:figure1}
% \end{figure}

% \begin{table}
%   \centering
%   \begin{tabular}{|c|c|c|}
%     \hline
%     \tabhead{Objects} &
%     \multicolumn{1}{|p{0.3\columnwidth}|}{\centering\tabhead{Caption --- pre-2002}} &
%     \multicolumn{1}{|p{0.4\columnwidth}|}{\centering\tabhead{Caption --- 2003 and afterwards}} \\
%     \hline
%     Tables & Above & Below \\
%     \hline
%     Figures & Below & Below \\
%     \hline
%   \end{tabular}
%   \caption{Table captions should be placed below the table.}
%   \label{tab:table1}
% \end{table}



% Balancing columns in a ref list is a bit of a pain because you
% either use a hack like flushend or balance, or manually insert
% a column break.  http://www.tex.ac.uk/cgi-bin/texfaq2html?label=balance
% multicols doesn't work because we're already in two-column mode,
% and flushend isn't awesome, so I choose balance.  See this
% for more info: http://cs.brown.edu/system/software/latex/doc/balance.pdf
%
% Note that in a perfect world balance wants to be in the first
% column of the last page.
%
% If balance doesn't work for you, you can remove that and
% hard-code a column break into the bbl file right before you
% submit:
%
% http://stackoverflow.com/questions/2149854/how-to-manually-equalize-columns-
% in-an-ieee-paper-if-using-bibtex
%
% Or, just remove \balance and give up on balancing the last page.
%
\newpage
\balance

% REFERENCES FORMAT
% References must be the same font size as other body text.

\bibliographystyle{acm-sigchi}
\bibliography{masterthesis-pape}


\clearpage
\newpage
\onecolumn

\begin{landscape}


\section{Appendix A. Semi-Structured Interview and Expert Interview Guidelines}
\subsubsection{Appendix A.1. Semi-Structured Interview Guideline (in German)}
The interview guideline was developed following rules of Helfferich \cite{helfferich2005}. It is in german as the interviews were held in german. Participants were shown the developed application prior to the interview.
\begin{adjustwidth}{-8em}{-8em}

\begin{longtable}{|p{6.45cm}|p{6.45cm}|p{6.45cm}|p{6.45cm}|}
 \hline
 \textbf{Leitfrage (Erz{\"a}hlaufforderung)}&\textbf{Check -- Wurde das erw{\"a}hnt? Memo f{\"u}r m{\"o}gliche Nachfragen -- nur stellen wenn nicht von allein angesprochen! Formulierung anpassen}&\textbf{Konkrete Fragen -- bitte an passender Stelle (auch am Ende m{\"o}glich) in dieser Formulierung stellen}&\textbf{Aufrechterhaltungs- und Steuerungsfragen}\\
 \hline

 \multicolumn{4}{|l|}{\textbf{Teil 1 -- B{\"u}rgerbeteiligung}}\\
 \hline
 
 Erz{\"a}hlen Sie mir {\"u}ber ihre Rolle und Aufgaben in B{\"u}rgerbeteiligung & Wie lange aktiv (Befragter, Projekt)\newline "`Organisator"' oder "`an der Basis"' & & Erz{\"a}hlen Sie noch mehr {\"u}ber\dots \\
 \hline
 
 
 Bitte beschreiben Sie mir die aus ihrer Sicht wichtigsten Aspekte von B{\"u}rgerbeteiligung. & Ziele \newline Nutzen & & \\
 \hline
 
 Bitte geben Sie mir eine Einf{\"u}hrung in ein(e) laufende(s)/ abgeschlossene(s) Initiative/Projekt (spontan entscheiden welches mehr "`dialogische"' Interaktion zwischen B{\"u}rgern und Aktion erfordert)& Methoden f{\"u}r B{\"u}rgerbefragung \newline Wie erfolgreich/Probleme? \newline "`Moderne"' (Social media) methoden angedacht? \newline Form von Beitr{\"a}gen die B{\"u}rger gebracht haben \newline Wie wurden die Aspekte ber{\"u}cksichtigt?
   & Welchen Wert wurde auf Dialoge zwischen den Akteuren gelegt? & Wie ist das ganze dann abgelaufen?\\
 \hline
 
 \multicolumn{4}{|l|}{\textbf{Teil 2 -- Einsatz der Anwendung}}\\
 \hline
 
 Bitte geben Sie mir eine Einf{\"u}hrung in das Projekt in dem Sie die Anwendung einsetzen wollen. & Zielgruppe (Bev{\"o}lkerungsgruppen, Geographisch) \newline redaktionelle Inhalte \newline erwartete Inhalte \newline Anreize zu Dialogen/Austausch mit B{\"u}rgern? & K{\"o}nnen Sie sich weitere Anwendungsf{\"a}lle f{\"u}r die Verkn{\"u}pfung von Texten mit Karten neben B{\"u}rgerbeteiligung vorstellen? & Erz{\"a}hlen Sie noch mehr {\"u}ber\dots \\
 \hline
 
 Welche Gr{\"u}nde sprechen f{\"u}r den Einsatz dieser L{\"o}sung gegen{\"u}ber anderen L{\"o}sungen. & Bedingungen (technisch, funktional) \newline angedachte Alternativen und deren Defizite \newline B{\"u}rgerbeteiligungsaspekte ber{\"u}cksichtigt? & Welche Eigenschaften w{\"u}rden Sie davon abhalten solch eine Anwendung einzusetzen? \newline Was k{\"o}nnte B{\"u}rger davon abhalten sich durch die Anwendung zu beteiligen? & \\
 \hline

 \multicolumn{4}{|l|}{\textbf{Teil 3 -- Abschlie{\ss}ende Fragen}}\\
 \hline
 
 Kennen Sie Beispiele f{\"u}r die Verkn{\"u}pfung geographischer Daten mit Diskussionsbeitr{\"a}gen? & Next Kassel/Hamburg \newline Frankfurt Gestalten \newline Shareabouts \newline collaborativemap.org & & \\
 \hline
 
 Haben Sie sich dort beteiligt? & In welcher Form & & Wie ist das ganze dann abgelaufen? \\
 \hline
 
 Kennen Sie Werkzeuge um interaktive Karten mit eigenen Inhalten zu erzeugen? & Google Map Maker \newline Here Map Creator \newline Wikimapia \newline Unclemap & & \\
 \hline
 
 Haben Sie schonmal ein solches Werkzeug eingesetzt? & Wie? & &\\
 \hline
 
 \end{longtable}
\end{adjustwidth}
\end{landscape}
\newpage

\subsubsection{Appendix A.1. Expert Interview Guideline}

As mentioned by Helfferich \cite{helfferich2005}, can handle more direct questions. Therefore these questions are much more straightforward. Because the interviews were held in german, the questions are also in german. Prior to the interview, the developed application was demoed.
\begin{itemize}
    \item Kennen Sie Anwendungen die Diskussionen durch Geoobjekte unterst{\"u}tzen?
    \item Welche davon haben Sie in der Vergangenheit schon einmal benutzt?
    \item Z{\"a}hlen Sie bitte die Vor- und Nachteile dieser Anwendungen auf
    \item Welche Anwendungsf{\"a}lle f{\"u}r die Verkn{\"u}pfung von Diskussionen und Geoobjekten k{\"o}nnen Sie sich außerhalb des B{\"u}rgerbeteiligungskontextes vorstellen?
    \item Welche L{\"o}sungen um B{\"u}rger mit Initiativen/Politik zusammenzubringen kennen Sie?
    \item Wie l{\"a}uft die Kommunikation zwischen den B{\"u}rgern und Initiativen/Politik bei diesen L{\"o}sungen ab?
    \item Denken Sie die explizite Verkn{\"u}pfung von Geoobjekten mit Diskussionsgegenst{\"a}nden ist generell hilfreich im B{\"u}rgerbeteiligungskontext / bei Dialogen?
    \item Im Vergleich zu den Anwendungen die Sie kennen, was denken Sie {\"u}ber die folgenden Funktionen der eben vorgestellten Anwendung?
        \begin{itemize}
            \item Verstecken von Geoobjekten die zu Antworten erstellt worden sind; In der Themenansicht nur die Geoobjekte der initialen Beitr{\"a}ge auf der Karte
            \item Zwei Wege Highlights von Geoobjekten und Beitr{\"a}gsboxen
            \item Filter und Sortierung
            \item Verfassen/Antworten
            \item Verkn{\"u}pfen von W{\"o}rtern mit neuen Geoobjekten, bestehenden Geoobjekten und Links
            \item Favorisierung von Beitr{\"a}gen
            \item Benutzerregistrierung/Anmeldung (und Social Login)
        \end{itemize}
    \item Werden ihrer Meinung nach Dialoge vereinfacht oder unterst{\"u}tzt?
    \item Welche Funktionen haben Sie vermisst?

\end{itemize}

\twocolumn
\section{Appendix B. Transcribed Interviews}

\subsubsection{Appendix B.1. Transcription System}
The interviews and the focus group were transcribed the following these rules (Rules from Kuckarz \cite{kuckartz2007} with modifications):
\begin{enumerate}
    \item The transcription is literal. Dialects are not transcribed.
    \item Punctuation and language are modified to match grammar and syntax of the german language.
    \item All personal details and mentions are removed and anonymized to prevent re-identification.
    \item Pauses and breaks are marked with ellipses (\dots).
    \item Agreeing sounds like ``Mhms'', ``Ahas'', etc. of the interviewer are not transcribed if they did not interrupt the interviewee.
    \item Interjections of the other person are in brackets.
    \item Supporting or clarifying sounds of the interviewee like laughing or sighing are noted in brackets.
    \item Passages of the interviewing person are denoted with ``I:'', passages of the interviewed person with a distinct abbreviation like ``P1:''.
\end{enumerate}

\subsubsection{Appendix B.2. Participant 1}

\textbf{Teil 1 -- B{\"u}rgerbeteiligung}
\begin{itemize}
    \item[I:] Erz{\"a}hlen Sie mir {\"u}ber ihre Rolle und Aufgaben in B{\"u}rgerbeteiligung
    \item[P1:] Ich kann nur sagen was f{\"u}r mich wichtig ist
    \item[I:] Ja das ist auch eine Sache die Sie mir erz{\"a}hlen k{\"o}nnen. Dann beschreiben Sie mir bitte die aus ihrer Sicht wichtigsten Aspekte der B{\"u}rgerbeteiligung
    \item[P1:] Die B{\"u}rgerbeteiligung f{\"u}hrt dazu, das erstmal Leute sich informieren, dass sie mehr wissen als nur {\"u}ber Zeitung. Dann k{\"o}nnen sie sich auch zusammenschlie{\ss}en und diskutieren und Aktionen besprechen. Ja und auch entsprechend Aktionen machen. Das st{\"a}rkt im Grunde eine Stadt auch.
    \item[I:] An welchen B{\"u}rgerbeteiligungsaktionen haben Sie dann schonmal teilgenommen?
	\item[P1:] Ja die Frage ist jetzt was alles unter B{\"u}rgerbeteiligung f{\"a}llt?
	\item[I:] Da kann alles drunter fallen, was f{\"u}r die Öffentlichkeit geschieht.
	\item[P1:] Also, ich habe zum Beispiel mit Leuten zusammen einen Gemeinschaftsgarten, der ist gegr{\"u}ndet worden und da treffen wir uns, da wird Gem{\"u}se angebaut, da gibts Bienen und da ist auch gedacht, dass man sich noch in Zukunft wenn das mal l{\"a}uft sich vernetzt mit anderen G{\"a}rten. Dass man vielleicht zu bestimmten Themen haben wir jetzt untereinander zum Beispiel zum Thema Bienen haben wir einen Nachmittag gehabt. Aber das kann man nat{\"u}rlich auch in einen gr{\"o}{\ss}eren Ma{\ss}stab machen.
	\item[I:] Und w{\"u}rden Sie denken, dass so eine Art von Anwendung dann in diesem Kontext sinnvoll w{\"a}re einzusetzen um das ganze bekannter zu machen und die Inhalte nach au{\ss}en zu kommunizieren?
	\item[P1:] Ja erstmal stell ich mir das so vor, dass jemand der, meinetwegen, neu ist oder keine Kontakte hat, sich mit Hilfe der Karte {\"u}berhaupt mal ein Bild machen was es f{\"u}r M{\"o}glichkeiten gibt. Und dann geht es ja in die Feindifferenzierung. Da w{\"u}rde er sagen: Gut, ich interessiere mich f{\"u}r Umwelt. Wer ist zust{\"a}ndig f{\"u}r Umwelt. Naja Greenpeace kann ich mal anklicken. Wo treffen die sich. Wann treffen die sich. Was haben die f{\"u}r Aktivit{\"a}ten zum Beispiel. Oder Transition Town. Was machen die eigentlich. Muss ich mal lesen was das {\"u}berhaupt ist. Ich wei{\ss} gar nicht genau was das ist. Also kann ich das mal lesen und vielleicht auch Kontakt aufnehmen."' Ich muss mir jetzt nicht m{\"u}hsam diese ganzen Adressen zusammensuchen. Diese WWW-Adressen, sondern die sind ja auf deiner Karte schon angegeben. Das ist nat{\"u}rlich schon auch erleichternd ist. Denn manchmal scheitert es an solchen Sachen. Auch an Bequemlichkeit.
	\item[I:] Wie l{\"a}uft dann im Moment die Kommunikation intern f{\"u}r diesen Garten ab?
	\item[P1:] Über E-Mail und {\"u}ber Treffen.
	\item[I:] Wie oft treffen sie sich da?
	\item[P1:] Ja da gibts dann Einladungen. Aber das ist unterschiedlich. Alle zwei Monate wenn was ansteht. Jetzt wo das Wasser da ist, da trifft man sich mal um aufzur{\"a}umen oder um Projekte zu besprechen.	
\end{itemize}

\textbf{Teil 2 -- Einsatz der Anwendung}
\begin{itemize}
	\item[I:] Wie soll dann die Beteiligung von Transition Town oder dem Garten auf dem Nachhaltigkeitstag aussehen?
	\item[P1:] Naja das kann ich jetzt nur erfinden. Letztenendes m{\"u}ssen wir das ja als Gruppe besprechen.
	\item[I:] Also am besten wie Sie sich das vorstellen
	\item[P1:] Themen die Transition Town wichtig sind w{\"u}rden da einen Raum finden und den Rahmen m{\"u}ssen die sich geben. Ob das jetzt in Form von (\dots) dass man gesundes Essen anbietet, oder mal so ne Karte entwirft wo Transition Town ist. Es gibt ja auch einen Film {\"u}ber Transition Town. Da gibts ja vielf{\"a}ltige M{\"o}glichkeiten. (\dots) Zum Beispiel in unserem Garten da hat die Bienen-Frau einen Vortrag gehalten {\"u}ber die Bienen und das soziale Miteinander. Das ist ja hoch differenziert. Zum Beispiel die Drohnen, die treffen sich an ganz bestimmten Pl{\"a}tzen vierzig Meter {\"u}ber der Erde. Solche die Detailinformationen die kein Mensch eigentlich wei{\ss}, die k{\"o}nnte man dann geben, in dem diese Bienen-Frau vielleicht was mitbringt und dann dar{\"u}ber redet und das dann auch Kindern zeigt wie so ein Bienenstock aussieht und mal Honig probieren l{\"a}sst. Und dann auch zum engagieren auffordert. Oder hab ich heute in der Zeitung gelesen, dass es einen Jungen gibt, der hat ein Bienenhaus gebaut f{\"u}r den Balkon. Der w{\"u}rde dann eingeladen und w{\"u}rde das vorstellen.
	\item[I:] Wer w{\"a}re dann die Zielgruppe? 
	\item[P1:] Wie meinst du die Zielgruppe?
	\item[I:] Ich meine damit die Personenkreise die man ansprechen m{\"o}chte
	\item[P1:] Naja an dem Tag werden ja viele Menschen da sein. Und das w{\"a}re ja dann ein wichtiger Aspekt zum Nachhaltigkeitsthema. Ich mein, da dass ja wahrscheinlich drau{\ss}en stattfindet, kann man ja von Zielgruppe nicht so unbedingt sprechen, oder? Wer will, der kommt.
	\item[I:] Gibt es andere Ans{\"a}tze die Sie zur Kommunikation bez{\"u}glich des Nachhaltigkeitstages in Betracht gezogen haben?
	\item[P1:] Nein im Moment nicht.
	\item[I:] Was f{\"u}r Gr{\"u}nde w{\"u}rden f{\"u}r den Einsatz der Karte sprechen? 
    \item[P1:] Also du meinst was f{\"u}r Vorteile es f{\"u}r uns h{\"a}tte den Garten in deine Karte einzutragen?
    \item[I:] Richtig.
	\item[P1:] Das h{\"a}tte den Vorteil, dass man auf einen Blick sehen kann, da und da und da gibt es einen freien Garten. Man sieht welche Adresse das sind. Man sieht vielleicht auch wann die da sind. Und dann ist das nat{\"u}rlich sehr {\"u}bersichtlich. Mit einem Klick hat man sozusagen die Information. Es gibt ja noch mehrere G{\"a}rten. Es gibt da unseren Paradies-Garten, dann gibt es am Campus noch einen Garten, dann gibts noch an der Gasselstiege einen Garten. Ja. Die w{\"u}rde man dann da sehen und dann k{\"o}nnte man auch Leute die das wollen, meinetwegen, eine Fahhradtour machen lassen und die G{\"a}rten angucken.
	\item[I:] Die Karte k{\"o}nnte man auch benutzen um die Fahrradtour zu organisieren. Dass man Start, Ziel und Zwischenhalte markiert. 
	\item[P1:] Ja genau.
	\item[I:] Was k{\"o}nnten Sachen sein die B{\"u}rger davon abhalten w{\"u}rden diese Karte zu benutzen?
	\item[P1:] (\dots) Ja also die Karte ist ja elektronisch. Geht ja nur {\"u}ber das Internet. Also sofern man einen Internetanschluss hat und einen Laptop oder einen Computer, gibts da nichts was dagegen spricht.
\end{itemize}
\textbf{Teil 3 -- Abschlie{\ss}ende Fragen}
\begin{itemize}
    \item[I:] Kennen Sie Beispiele f{\"u}r die Verkn{\"u}pfung geographischer Daten mit Diskussionsbeitr{\"a}gen?
	\item[P1:] Sag mir nochmal was man alles unter geographische Daten fasst.
	\item[I:] Orte und Objekte mit einem Ort
	\item[P1:] Naja ich war jetzt auf dem Jakobsweg, da hat man auch Karten. Aber die nutzt man nicht so oft. Da hat man B{\"u}cher in denen die Adressen drin stehen. Und die Zeichen sind an den B{\"a}umen.
	\item[I:] Ja ist auch ne M{\"o}glichkeit. Ich ziele mit der Frage eher ab auf auf elektronische Anwendungen.
	\item[P1:] Nein, ich nicht da auch nicht so firm. Ich mag das auch nicht.
	\item[I:] Also haben Sie sowas auch noch nie benutzt?
	\item[P1:] Nein.
	\item[I:] Kennen Sie Werkzeuge um interaktive Karten mit eigenen Inhalten zu erzeugen?
	\item[P1:] Nein. Es gibt ja viele Leute die nicht so interessiert sind mit den neuen Medien. 
	\item[I:] Ja. Alles klar. Das waren dann die Fragen von meiner Seite. Gibt es noch Fragen von ihrer Seite? 
    \item[P1:] Nein. Eigentlich nicht. Gute Sache.
    \item[I:] Vielen Dank f{\"u}r ihre Zeit.
\end{itemize}

\end{document}
